\chapter{Experiments, Sample Spaces and Counting}
\label{ch:counting}

Probability is a language for describing uncertain outcomes. Before we can
assign a number to ``how likely'', we must say precisely \emph{what can
happen}. This chapter builds that vocabulary (sample spaces and events), states
the classical rule for equally likely outcomes, and develops the counting tools
needed to apply it.

\section{Random experiments and sample spaces}

A \emph{random experiment} is a procedure whose result is not determined in
advance, but whose set of possible results is known. We do not need the
procedure to be repeatable in practice; we only need a clear description of the
possible results.

\begin{definition}[Sample space]
The \emph{sample space} $\Omega$ is the set of all possible outcomes of a
random experiment. An element $\omega\in\Omega$ is called an \emph{outcome}.
\end{definition}

The sample space must satisfy two requirements. First, it must be
\emph{exhaustive}: every possible result corresponds to some $\omega$. Second,
it must be \emph{mutually exclusive}: two different outcomes cannot occur
together in a single run of the experiment.

\begin{example}\label{ex:spaces}
\begin{enumerate}[label=(\alph*)]
\item One coin toss: $\Omega=\{H,T\}$.
\item One roll of a die: $\Omega=\{1,2,3,4,5,6\}$.
\item Two rolls of a die, recording both numbers in order:
$\Omega=\{1,\dots,6\}^2=\{(i,j):1\le i,j\le 6\}$, with $36$ outcomes.
\item Tossing a coin until the first head: $\Omega=\{1,2,3,\dots\}=\N$, where
$n$ means ``the first head occurs on toss $n$''. This set is infinite but
\emph{countable}.
\item The lifetime of a light bulb in hours: $\Omega=[0,\infty)$. This set is
\emph{uncountable}.
\end{enumerate}
\end{example}

\begin{nuance}{the sample space is a modelling choice}
The same physical experiment can be described by different sample spaces. For
two dice we could record the ordered pair $(i,j)$, or only the sum
$\{2,\dots,12\}$. Both are valid sample spaces, but they are not equally
convenient. A finer description (the ordered pair) always contains the coarser
one (the sum is a function of the pair). As we will see in
Section~\ref{sec:classical}, choosing a sample space whose outcomes are
equally likely is often the key step in a solution.
\end{nuance}

\section{Events and the language of sets}

We are rarely interested in a single outcome. We want to ask questions such as
``is the sum at least 10?'' Each such question corresponds to a subset of
$\Omega$.

\begin{definition}[Event]
An \emph{event} is a subset $A\subseteq\Omega$. We say that $A$ \emph{occurs}
if the outcome $\omega$ of the experiment satisfies $\omega\in A$.
\end{definition}

For a finite or countably infinite $\Omega$ we allow every subset to be an
event. For uncountable $\Omega$ this is not always possible; the precise
resolution (a $\sigma$-algebra of events) is given in
Chapter~\ref{ch:axioms}, and the reason it is necessary is explained in
Chapter~\ref{ch:measure}. Until then, ``event'' means ``any subset we can describe''.

Combining events corresponds to combining sets.

\begin{center}
\begin{tabular}{@{}lll@{}}
\toprule
In words & Notation & Meaning \\
\midrule
$A$ does not occur & $A^c$ & $\{\omega\in\Omega:\omega\notin A\}$\\
$A$ or $B$ (or both) occurs & $A\cup B$ & $\omega\in A$ or $\omega\in B$\\
$A$ and $B$ both occur & $A\cap B$ & $\omega\in A$ and $\omega\in B$\\
$A$ occurs but $B$ does not & $A\setminus B=A\cap B^c$ & \\
$A$ and $B$ cannot both occur & $A\cap B=\emptyset$ & \emph{disjoint} or \emph{mutually exclusive}\\
if $A$ occurs then $B$ occurs & $A\subseteq B$ & \\
\bottomrule
\end{tabular}
\end{center}

The event $\Omega$ always occurs (the \emph{certain} event) and $\emptyset$ never
occurs (the \emph{impossible} event).

\begin{theorem}[De Morgan's laws]\label{thm:demorgan}
For any events $A_1,A_2,\dots$ (finitely or countably many),
\[
\Bigl(\bigcup_i A_i\Bigr)^c=\bigcap_i A_i^c,
\qquad
\Bigl(\bigcap_i A_i\Bigr)^c=\bigcup_i A_i^c .
\]
\end{theorem}
\begin{proof}
We prove the first identity by showing each side is contained in the other. Let
$\omega\in(\bigcup_i A_i)^c$. Then $\omega$ lies in no $A_i$, so
$\omega\in A_i^c$ for every $i$, that is, $\omega\in\bigcap_i A_i^c$.
Conversely, if $\omega\in\bigcap_i A_i^c$, then $\omega\notin A_i$ for every
$i$, so $\omega\notin\bigcup_i A_i$. The second identity follows by applying the
first to the events $A_i^c$ and taking complements.
\end{proof}

\begin{definition}[Partition]
A finite or countable collection of events $B_1,B_2,\dots$ is a
\emph{partition} of $\Omega$ if the events are pairwise disjoint
($B_i\cap B_j=\emptyset$ for $i\ne j$) and $\bigcup_i B_i=\Omega$. Exactly one
$B_i$ occurs on every run of the experiment.
\end{definition}

\section{The classical definition of probability}\label{sec:classical}

Historically, probability began with games of chance in which, by symmetry,
every outcome was judged equally likely.

\begin{definition}[Classical probability]
Let $\Omega$ be a \emph{finite} sample space whose outcomes are equally likely.
The probability of an event $A\subseteq\Omega$ is
\[
\Prob(A)=\frac{\abs{A}}{\abs{\Omega}} ,
\]
where $\abs{\cdot}$ denotes the number of elements.
\end{definition}

Everything depends on the phrase ``equally likely''. It is an \emph{assumption
about the model}, justified (or not) by physical symmetry. It is not something
the formula can verify.

\begin{example}[Two dice]
Take $\Omega=\{1,\dots,6\}^2$, which has $36$ equally likely outcomes. The event
``sum equals $7$'' is $\{(1,6),(2,5),(3,4),(4,3),(5,2),(6,1)\}$, so its
probability is $6/36=1/6$.
\end{example}

\begin{pitfall}{equal likelihood on the wrong space}
Suppose we toss two fair coins and let $N$ be the number of heads. It is
tempting to take $\Omega'=\{0,1,2\}$ and conclude $\Prob(N=1)=1/3$. This is
wrong. The three values of $N$ are not equally likely. The four ordered
outcomes $HH,HT,TH,TT$ are equally likely, and $N=1$ corresponds to two of
them, so $\Prob(N=1)=2/4=1/2$. A famous historical error of exactly this kind is
attributed to d'Alembert, who counted two coin tosses as having three
equally likely outcomes. The same trap arises when summing two dice and
treating the eleven possible sums as equally likely. The remedy: \emph{when physical symmetry is
the justification, apply it to the finest description in which the objects
are distinguishable}.
\end{pitfall}

\begin{nuance}{limits of the classical definition}
The formula $\abs{A}/\abs{\Omega}$ fails when $\Omega$ is infinite (there is no
way to choose a natural number ``uniformly''; see
Chapter~\ref{ch:axioms}), and it is circular as a \emph{definition}, since
``equally likely'' already means ``equally probable''. In a continuous
setting, ``choose at random'' can be ambiguous: Bertrand's paradox asks for the
probability that a random chord of a circle is longer than the side of the
inscribed equilateral triangle, and three natural interpretations of ``random
chord'' give $1/3$, $1/2$ and $1/4$. The paradox shows that the sampling
mechanism must be specified. Chapter~\ref{ch:axioms} gives the axiomatic
framework that removes the circularity: probability becomes a function
satisfying rules, and the classical formula becomes a special case.
\end{nuance}

\section{Counting}\label{sec:counting}

With equally likely outcomes, probability reduces to counting. This section
develops the standard tools with proofs.

\subsection{Basic principles}

\begin{proposition}[Rule of sum and rule of product]\label{prop:rules}
\begin{enumerate}[label=(\roman*)]
\item \emph{Sum rule.} If a task can be done in one of $m$ ways or in one of
$n$ ways, and no way is counted twice, there are $m+n$ ways.
\item \emph{Product rule.} If a procedure has $k$ stages, and stage $i$ can be
carried out in $n_i$ ways regardless of the choices made in earlier stages,
then the whole procedure can be carried out in $n_1n_2\cdots n_k$ ways.
\end{enumerate}
\end{proposition}
\begin{proof}
The sum rule is the statement $\abs{A\cup B}=\abs{A}+\abs{B}$ for disjoint
finite sets. For the product rule with two stages, list the choices of stage~1
as $a_1,\dots,a_{n_1}$. For each $a_i$ there are $n_2$ completions, giving
$n_2$ distinct pairs $(a_i,\cdot)$. These blocks are disjoint, so the total is
$n_1n_2$. Induction on $k$ gives the general case.
\end{proof}

\begin{nuance}{``regardless of earlier choices'' means the \emph{number}, not the set}
In stage 2 the available options may depend on stage 1, as long as their
\emph{count} does not. For example, choosing a president and then a vice
president from $n$ people gives $n(n-1)$ outcomes: the set of candidates for
vice president depends on who is president, but there are always $n-1$ of
them.
\end{nuance}

\subsection{Permutations and combinations}

\begin{definition}[Factorial]
For an integer $n\ge 0$, $n!=n(n-1)\cdots 2\cdot 1$, with $0!=1$.
\end{definition}

\begin{proposition}[Arrangements]\label{prop:perm}
\begin{enumerate}[label=(\roman*)]
\item The number of ways to arrange $n$ distinct objects in a row is $n!$.
\item The number of ordered selections of $k$ objects from $n$ distinct
objects, without repetition (called $k$-\emph{permutations}), is
\[
P(n,k)=n(n-1)\cdots(n-k+1)=\frac{n!}{(n-k)!}.
\]
\end{enumerate}
\end{proposition}
\begin{proof}
For (ii), apply the product rule: $n$ choices for the first position, $n-1$ for
the second (whatever the first was), and so on, down to $n-k+1$ for the $k$th.
Part (i) is the case $k=n$.
\end{proof}

\begin{definition}[Binomial coefficient]
The number of $k$-element subsets of an $n$-element set is denoted
$\binom{n}{k}$ (read ``$n$ choose $k$''), for integers $0\le k\le n$. We set
$\binom{n}{k}=0$ if $k<0$ or $k>n$.
\end{definition}

\begin{theorem}[Formula for $\binom nk$]\label{thm:choose}
For $0\le k\le n$,
\[
\binom{n}{k}=\frac{n!}{k!\,(n-k)!}.
\]
\end{theorem}
\begin{proof}
Count the ordered selections of $k$ distinct objects from $n$ in two ways. By
Proposition~\ref{prop:perm} there are $n!/(n-k)!$ of them. Alternatively, first
choose which $k$-element \emph{set} is selected ($\binom nk$ ways), then order
its elements ($k!$ ways, by Proposition~\ref{prop:perm}(i)). By the product
rule the count is $\binom nk\,k!$. Equating the two counts and dividing by
$k!$ gives the formula.
\end{proof}

Two immediate consequences: \emph{symmetry}, $\binom nk=\binom n{n-k}$
(choosing the $k$ objects to keep is the same as choosing the $n-k$ to discard),
and the special values $\binom n0=\binom nn=1$, $\binom n1=n$.

\begin{theorem}[Pascal's rule]\label{thm:pascal}
For $1\le k\le n-1$,
\[
\binom nk=\binom{n-1}{k-1}+\binom{n-1}{k}.
\]
\end{theorem}
\begin{proof}
Fix one special object $x$ among the $n$. Every $k$-subset either contains $x$
or does not. Those containing $x$ are formed by choosing the other $k-1$
members from the remaining $n-1$ objects: $\binom{n-1}{k-1}$ ways. Those not
containing $x$ are $k$-subsets of the remaining $n-1$ objects:
$\binom{n-1}{k}$ ways. The two cases are disjoint and exhaustive; apply the sum
rule.
\end{proof}

\begin{theorem}[Binomial theorem]\label{thm:binomial}
For real $a,b$ and integer $n\ge 0$,
\[
(a+b)^n=\sum_{k=0}^n\binom nk a^k b^{n-k}.
\]
\end{theorem}
\begin{proof}
Write $(a+b)^n=(a+b)(a+b)\cdots(a+b)$ with $n$ factors. Expanding, we choose
from each factor either $a$ or $b$, obtaining $2^n$ terms. The term $a^kb^{n-k}$
arises exactly when $a$ is chosen from some set of $k$ of the $n$ factors; there
are $\binom nk$ such sets. Collecting terms gives the formula.
\end{proof}

\begin{corollary}\label{cor:sums}
For $n\ge 1$: $\displaystyle\sum_{k=0}^n\binom nk=2^n$ and
$\displaystyle\sum_{k=0}^n(-1)^k\binom nk=0$.
\end{corollary}
\begin{proof}
Put $(a,b)=(1,1)$ and $(a,b)=(-1,1)$ in Theorem~\ref{thm:binomial}.
\end{proof}

The first identity also has a direct meaning: a set with $n$ elements has $2^n$
subsets, since each element is independently in or out.

\begin{theorem}[Vandermonde's identity]\label{thm:vander}
For non-negative integers $m,n,r$,
\[
\sum_{k=0}^r\binom mk\binom n{r-k}=\binom{m+n}{r}.
\]
\end{theorem}
\begin{proof}
Take a group of $m$ men and $n$ women and count the committees of size $r$. The
right side counts all committees. On the left, $k$ is the number of men on the
committee: choose them in $\binom mk$ ways and then the $r-k$ women in
$\binom n{r-k}$ ways. Summing over $k$ counts every committee once.
\end{proof}

\subsection{Arrangements with repeated objects}

\begin{theorem}[Multinomial coefficient]\label{thm:multinomial}
Suppose $n=n_1+\dots+n_r$ objects are arranged in a row, where $n_i$ of them are
of type $i$, and objects of the same type are indistinguishable. The number of
distinguishable arrangements is
\[
\binom{n}{n_1,\dots,n_r}=\frac{n!}{n_1!\,n_2!\cdots n_r!}.
\]
\end{theorem}
\begin{proof}
Give each object a unique label, so that there are $n!$ arrangements. Erasing
the labels groups arrangements together: two labelled arrangements look the
same exactly when one is obtained from the other by permuting labels within
each type. For each unlabelled arrangement there are $n_1!\cdots n_r!$ such
labelled arrangements (independent relabelling within each type, by the product
rule). Hence the number of unlabelled arrangements is
$n!/(n_1!\cdots n_r!)$.
\end{proof}

For example, the letters of MISSISSIPPI ($1$ M, $4$ I, $4$ S, $2$ P) can be
arranged in $\frac{11!}{1!\,4!\,4!\,2!}=34{,}650$ distinguishable ways.

\subsection{The four sampling schemes}

Choosing $k$ items from $n$ distinct items can be done with or without
replacement, and the result can be recorded as an ordered list or an unordered
collection. The four counts are:

\begin{center}
\begin{tabular}{@{}lcc@{}}
\toprule
 & Ordered & Unordered \\
\midrule
With replacement & $n^k$ & $\displaystyle\binom{n+k-1}{k}$ \\[0.6em]
Without replacement & $\displaystyle\frac{n!}{(n-k)!}$ & $\displaystyle\binom nk$ \\
\bottomrule
\end{tabular}
\end{center}

Three of these were derived above ($n^k$ is the product rule with $n$ options at
each of $k$ stages). The remaining one, unordered with replacement, is
equivalent to counting the number of multisets of size $k$ from $n$ types, and
needs a separate argument.

\begin{theorem}[Stars and bars]\label{thm:starsbars}
The number of non-negative integer solutions $(x_1,\dots,x_n)$ of
$x_1+\dots+x_n=k$ is $\binom{n+k-1}{k}$. Equivalently, this is the number of
multisets of size $k$ drawn from $n$ types.
\end{theorem}
\begin{proof}
A solution says how many items $x_i$ of type $i$ are chosen. Encode it as a row
of $k$ stars and $n-1$ bars: write $x_1$ stars, then a bar, then $x_2$ stars,
then a bar, and so on, ending with $x_n$ stars. For example, with $n=3,k=4$ the
solution $(2,0,2)$ is $\star\star\,|\,|\,\star\star$. This is a bijection between
solutions and rows with $k+n-1$ positions, of which exactly $k$ hold stars. A
row is determined by the set of positions of the stars, so there are
$\binom{n+k-1}{k}$ of them.
\end{proof}

\begin{nuance}{unordered samples are generally \emph{not} equally likely}
When items are drawn one at a time \emph{with} replacement from $n$ types, the
$n^k$ ordered sequences are equally likely, but the $\binom{n+k-1}{k}$ unordered
multisets are not. For $n=3$ types and $k=2$ draws, the multiset $\{a,a\}$
arises from one sequence $(a,a)$, whereas $\{a,b\}$ arises from two sequences,
$(a,b)$ and $(b,a)$. So $\Prob(\{a,a\})=1/9$ but $\Prob(\{a,b\})=2/9$. Applying
the classical formula to multisets would be wrong here. (Interestingly, there
are physical systems, namely bosons, for which the multisets \emph{are}
equally likely; whether objects are distinguishable is a fact about nature,
not a convention of mathematics.)
\end{nuance}

\section{Worked examples}

\begin{example}[Poker: a full house]\label{ex:fullhouse}
A poker hand is a $5$-card subset of a $52$-card deck; there are
$\binom{52}5=2{,}598{,}960$ equally likely hands. A \emph{full house} has three
cards of one rank and two of another. Count by stages: choose the rank of the
triple ($13$ ways), choose which $3$ of its $4$ suits ($\binom43=4$), choose a
different rank for the pair ($12$ ways), choose its $2$ suits
($\binom42=6$). Thus
\[
\Prob(\text{full house})=\frac{13\cdot 4\cdot 12\cdot 6}{2{,}598{,}960}
=\frac{3744}{2{,}598{,}960}\approx 0.00144 .
\]
\end{example}

\begin{example}[Exactly two aces]\label{ex:aces}
The probability that a $5$-card hand contains exactly two aces: choose the aces
($\binom42=6$) and the other three cards from the $48$ non-aces
($\binom{48}3=17{,}296$):
\[
\frac{6\cdot 17{,}296}{2{,}598{,}960}=\frac{103{,}776}{2{,}598{,}960}\approx 0.0399.
\]
This is a \emph{hypergeometric} probability; we will study that distribution in
Chapter~\ref{ch:dist}.
\end{example}

\begin{example}[The birthday problem]\label{ex:birthday}
In a room of $n$ people, what is the probability that at least two share a
birthday? Assume $365$ equally likely birthdays, independent across people
(ignoring leap years and seasonal effects). Take
$\Omega=\{1,\dots,365\}^n$, which has $365^n$ equally likely outcomes. It is
easier to count the complement $N$, ``all birthdays distinct'': by the product
rule there are $365\cdot364\cdots(365-n+1)$ such outcomes, so
\[
\Prob(N)=\prod_{k=0}^{n-1}\Bigl(1-\frac k{365}\Bigr),
\qquad
\Prob(\text{some match})=1-\Prob(N).
\]
To see why a small group suffices, use the inequality $1-x\le e^{-x}$ (valid
for all real $x$; it holds because $e^{-x}$ is convex and $1-x$ is its tangent
line at $x=0$). Applying it to each factor,
\[
\Prob(N)\le\prod_{k=0}^{n-1}e^{-k/365}=\exp\!\Bigl(-\frac{n(n-1)}{2\cdot 365}\Bigr),
\]
using $\sum_{k=0}^{n-1}k=n(n-1)/2$. For $n=23$ the exponent is $-506/730$,
which is slightly below $-\ln 2$ (because $730\ln 2\approx 505.98<506$). So
$\Prob(N)<1/2$ and $\Prob(\text{some match})>1/2$ when $n=23$. The exact value
is $\approx 0.507$.

The surprise comes from the number of \emph{pairs}: there are
$\binom{23}2=253$ pairs of people, and each pair matches with probability
$1/365$. The expected number of matching pairs is about $253/365\approx 0.69$,
which is not small.
\end{example}

\begin{interview}{count the complement, count pairs}
Two habits solve a large fraction of elementary counting questions. First,
``at least one'' events are usually easiest through the complement
``none''. Second, when many weakly dependent coincidences are possible, count
the number of \emph{pairs} (or tuples) that could produce one; this tells you
the right order of magnitude immediately. Both ideas will be formalised later:
the first through the axioms (Chapter~\ref{ch:axioms}), the second through
linearity of expectation (Chapter~\ref{ch:drv}).
\end{interview}

\section*{Exercises}
\addcontentsline{toc}{section}{Exercises}

\begin{enumerate}
\item A fair die is rolled twice. List the sample space and find the
probability that the second roll is strictly larger than the first.
\item Prove that $\binom nk=\frac nk\binom{n-1}{k-1}$ for $1\le k\le n$, once
algebraically and once by a counting argument (committee with a designated
chair).
\item Prove $\sum_{k=0}^n\binom nk^2=\binom{2n}{n}$ using Vandermonde's
identity.
\item How many ways are there to seat $n$ people at a round table, if seatings
that differ only by rotation are considered the same? Justify your answer.
\item How many integer solutions satisfy $x_1+x_2+x_3=10$ with every
$x_i\ge 2$? (Hint: shift the variables, then use stars and bars.)
\item Find the probability that a $5$-card hand is a flush (all one suit,
including straight flushes).
\item A password has $8$ characters, each a lowercase letter or digit. Find the
probability that a uniformly random password has no repeated character.
\item In the birthday problem, show that the exact probability of a match is
increasing in $n$. Then compute, using the bound in
Example~\ref{ex:birthday}, a value of $n$ guaranteeing a match probability of
at least $0.9$.
\item (Challenge) Prove the multinomial theorem
$(x_1+\dots+x_r)^n=\sum\binom{n}{n_1,\dots,n_r}x_1^{n_1}\cdots x_r^{n_r}$,
where the sum is over all non-negative integers with $n_1+\dots+n_r=n$.
\end{enumerate}
